\documentclass[11pt]{article}
\usepackage[margin=1in]{geometry}
\usepackage{amsmath,amssymb,amsthm,mathtools}
\usepackage{hyperref}
\usepackage{graphicx}
\usepackage{enumitem}

\newtheorem{theorem}{Theorem}
\newtheorem{lemma}{Lemma}
\newtheorem{remark}{Remark}

\DeclareMathOperator{\diag}{diag}
\DeclareMathOperator{\tr}{tr}
\newcommand{\C}{\mathbb{C}}
\newcommand{\one}{\mathbf{1}}

\title{A Similarity Proof of Conjecture 6.3 in\\
\emph{Sharp Schoenberg Type Inequalities and the de Bruin--Sharma Problem}}
\author{agent\_lane\_16}
\date{2026-07-18}

\begin{document}
\maketitle

\section*{Status}

\textbf{full\_solution\_likely\_valid.}
This packet proves Conjecture 6.3 of Tang--Zhang,
\emph{Sharp Schoenberg type inequalities and the de Bruin--Sharma problem},
arXiv:2508.10341v3.  The source verifies the conjecture only for \(n=2\);
the argument below proves it for every \(n\), and gives the equality
characterization.

\section{Source Problem}

The source considers normalized polynomials
\[
  P(z)=z\prod_{j=1}^n(z-z_j), \qquad z_j\in\C\setminus\{0\},
\]
and denotes by \(w_1,\ldots,w_n\) the zeros of \(P'(z)\).  Conjecture 6.3
states that
\begin{equation}\label{eq:C1}
  \sum_{k=1}^n\frac1{|w_k|^2}
  \le
  3\sum_{j=1}^n\frac1{|z_j|^2}
  +\left|\sum_{j=1}^n\frac1{z_j}\right|^2,
\end{equation}
and that equality holds whenever the points \(z_1,\ldots,z_n\) all lie on a
single line through the origin.

\begin{center}
\includegraphics[width=0.86\linewidth]{figures/open_problem_crop.png}
\end{center}

\section{Idea of the Proof}

The source already gives the decisive matrix model.  If
\[
  u_j=\frac1{z_j},\qquad u=(u_1,\ldots,u_n)^T,
\]
then \(1/w_1,\ldots,1/w_n\) are the eigenvalues of
\[
  B=\diag(u)+u\one^T=\diag(u)(I+J),
\]
where \(J=\one\one^T\).  The source estimates these eigenvalues by the
Frobenius norm of \(B\), which gives the coarser constant \(n+3\).

The small twist is to keep the same eigenvalues but change the Frobenius norm.
Let \(C=I+J\) and \(H=C^{1/2}\).  Since \(C\) is positive definite,
\[
  B=\diag(u)C
\]
is similar to
\[
  A=H\diag(u)H.
\]
Schur's inequality bounds the squared moduli of the eigenvalues by
\(\|A\|_F^2\), and \(\|A\|_F^2\) is exactly the right-hand side of
\eqref{eq:C1}.  Equality is Schur equality, i.e. normality of \(A\), and this
normality condition says precisely that all \(u_j\), equivalently all \(z_j\),
are collinear with the origin.

\section{Proof}

We use the following standard form of Schur's inequality.

\begin{lemma}[Schur inequality]\label{lem:schur}
Let \(A\in M_n(\C)\) have eigenvalues \(\lambda_1,\ldots,\lambda_n\). Then
\[
  \sum_{k=1}^n |\lambda_k|^2 \le \tr(AA^*)=\|A\|_F^2.
\]
Equality holds if and only if \(A\) is normal.
\end{lemma}

\begin{theorem}\label{thm:c1}
Let
\[
  P(z)=z\prod_{j=1}^n(z-z_j),\qquad z_j\in\C\setminus\{0\},
\]
and let \(w_1,\ldots,w_n\) be the zeros of \(P'(z)\). Then
\[
  \sum_{k=1}^n\frac1{|w_k|^2}
  \le
  3\sum_{j=1}^n\frac1{|z_j|^2}
  +\left|\sum_{j=1}^n\frac1{z_j}\right|^2.
\]
Moreover, equality holds if and only if \(z_1,\ldots,z_n\) lie on one line
through the origin.
\end{theorem}

\begin{proof}
Set \(u_j=1/z_j\), \(u=(u_1,\ldots,u_n)^T\), \(D=\diag(u)\), and
\(J=\one\one^T\).  The matrix model used in the source's proof of Theorem 6.1
shows that the reciprocal critical points \(1/w_1,\ldots,1/w_n\) are the
eigenvalues of
\[
  B=D+u\one^T = D(I+J).
\]
Put \(C=I+J\).  The matrix \(C\) is positive definite, so \(H=C^{1/2}\) is
invertible.  The matrix
\[
  A=HDH
\]
is similar to \(B=DC\), because
\[
  H^{-1}AH = H^{-1}(HDH)H = DC=B.
\]
Thus \(A\) and \(B\) have the same eigenvalues.  By Lemma \ref{lem:schur},
\[
  \sum_{k=1}^n\frac1{|w_k|^2}
  =\sum_{k=1}^n|\lambda_k(A)|^2
  \le \tr(AA^*).
\]
We now compute the last trace:
\[
  \tr(AA^*)=\tr(HDH^2D^*H)=\tr(CDCD^*).
\]
Since \(C=I+J\),
\[
  \tr(CDCD^*)
  =
  \tr(DD^*)+\tr(JDD^*)+\tr(DJD^*)+\tr(JDJD^*).
\]
The first three terms are all \(\sum_j |u_j|^2\).  For the last term,
\[
  \tr(JDJD^*)
  =(\one^T D\one)(\one^T D^*\one)
  =\left|\sum_{j=1}^n u_j\right|^2.
\]
Therefore
\[
  \tr(AA^*)
  =
  3\sum_{j=1}^n |u_j|^2
  +\left|\sum_{j=1}^n u_j\right|^2,
\]
which is exactly \eqref{eq:C1} after substituting \(u_j=1/z_j\).

It remains to discuss equality.  By Lemma \ref{lem:schur}, equality in
\eqref{eq:C1} holds if and only if \(A\) is normal.  Since \(H\) is positive
and invertible,
\[
  AA^*=A^*A
  \quad\Longleftrightarrow\quad
  DCD^*=D^*CD.
\]
The diagonal terms \(DD^*=D^*D\) cancel, so this is equivalent to
\[
  DJD^*=D^*JD.
\]
Entrywise, this says
\[
  u_i\overline{u_j}=\overline{u_i}u_j
  \qquad (1\le i,j\le n).
\]
Since no \(u_j\) is zero, the last condition is equivalent to every quotient
\(u_i/u_j\) being real, i.e. to the points \(u_1,\ldots,u_n\) lying on a single
real line through the origin.  Taking reciprocals preserves this property, so
this is equivalent to \(z_1,\ldots,z_n\) lying on a single line through the
origin.
\end{proof}

\section{Verification Notes}

The proof uses only:
\begin{itemize}[leftmargin=2em]
  \item the source paper's reciprocal-critical-point matrix model
  \(B=\diag(u)+u\one^T\);
  \item similarity by the positive square root of \(I+J\);
  \item Schur's eigenvalue-singular-value inequality.
\end{itemize}
The script \texttt{code/verify\_c1\_random.py} performs a randomized numerical
check for the equivalent matrix inequality.

\section{Novelty Check}

Local index searches found no exact packet for arXiv:2508.10341 or for
Conjecture 6.3.  Web searches on 2026-07-18 for exact phrases around
``Conjecture C1'', ``Conjecture 6.3'', the reciprocal critical-point
inequality, and the matrix forms \(\diag(u)+u\one^T\) and \(D^{-1}+b\one^T\)
found the source arXiv/J-GLOBAL records and adjacent Schoenberg-type material,
but no exact solution or this similarity argument.  Novelty remains subject to
human literature review.

\section{References}

\begin{enumerate}[leftmargin=2em]
  \item Q. Tang and T. Zhang,
  \emph{Sharp Schoenberg type inequalities and the de Bruin--Sharma problem},
  arXiv:2508.10341v3, 2025.
  \item I. Schur, standard eigenvalue majorization inequality; also used in
  the source paper as the eigenvalue--singular-value comparison lemma.
\end{enumerate}

\section{Human Review Recommendation}

Review as a likely full solution to Conjecture 6.3.  The key points to verify
are the source matrix model for reciprocal critical points, the similarity
\(D(I+J)\sim (I+J)^{1/2}D(I+J)^{1/2}\), and the equality characterization via
normality.

\end{document}
